\documentclass[../../main.tex]{subfiles}

\begin{document}

\section{引言}\label{section:ch12-引言}

本章研究半线性波动方程初值问题解的存在性（有时是不存在性）：
\begin{align}
u_{tt}-\Delta u+f(u)=0.
\label{align:ch12-半线性波动}
\end{align}
线性情形 $f(u)=m^2u$ 即 Klein--Gordon 方程。我们也将讨论某些具形式
\begin{align}
u_{tt}-\Delta u+f(Du,u_t,u)=0
\label{align:ch12-拟线性波动}
\end{align}
的温和拟线性波动方程，其中如通常写 $Du=D_xu$ 表示 $x$ 变量中的梯度。我们遵循习惯将非线性项放在等号左边，这简化了后面一些公式。更复杂的拟线性波动方程（其中二阶导数系数依赖于 $Du,u_t,u$）超出本书范围。

\subsection{能量守恒}\label{subsection:ch12-能量守恒}

首先考虑半线性 PDE \eqref{align:ch12-半线性波动}，设
\begin{align}
F(z)\coloneqq\int_0^z f(w)\,\mathrm{d}w \quad (z\in\mathbb{R}).
\label{align:ch12-F定义}
\end{align}
则 $F(0)=0$ 且 $F'=f$。回忆 \S8.6.2，解 $u$ 在时刻 $t\geqslant 0$ 的能量为
\begin{align}
E(t)\coloneqq\int_{\mathbb{R}^n}\frac{1}{2}(u_t^2+|Du|^2)+F(u)\,\mathrm{d}x.
\label{align:ch12-能量}
\end{align}

\begin{theorem}[能量守恒。]\label{theorem:ch12-能量守恒}假设 $u$ 是半线性波动方程\eqref{align:ch12-半线性波动} 的光滑解，且对每个时刻 $t$，$u(\cdot,t)$ 具有空间紧支集。则映射 $t\mapsto E(t)$ 为常数。
\end{theorem}

\begin{proof}
计算
\begin{align*}
\dot{E}(t)=\int_{\mathbb{R}^n}u_tu_{tt}+Du\cdot Du_t+f(u)u_t\,\mathrm{d}x=\int_{\mathbb{R}^n}u_t(u_{tt}-\Delta u+f(u))\,\mathrm{d}x=0,
\end{align*}
其中 $\cdot=\dfrac{\mathrm{d}}{\mathrm{d}t}$。
\end{proof}

此证明中的分部积分是有效的，因为 $u$ 对每个时刻具有空间紧支集。在后面许多证明中我们将类似地分部积分，隐式依赖于解在 $|x|$ 大时消失以证明计算的合理性。

\subsection{有限传播速度}\label{subsection:ch12-有限传播速度}

回忆 \S2.4.3 中线性波动方程的依赖域计算，重新引入后向波锥：

\begin{definition}\label{definition:ch12-后向波锥}
固定 $x_0\in\mathbb{R}^n$，$t_0>0$，定义顶点在 $(x_0,t_0)$ 的\textbf{后向波锥}：
\begin{align*}
K(x_0,t_0)\coloneqq\{(x,t)\mid 0\leqslant t\leqslant t_0,\;|x-x_0|\leqslant t_0-t\}.
\end{align*}
$K(x_0,t_0)$ 边界的弯曲部分为
\begin{align*}
\Gamma(x_0,t_0)\coloneqq\{(x,t)\mid 0\leqslant t\leqslant t_0,\;|x-x_0|=t_0-t\}.
\end{align*}
\end{definition}

\begin{theorem}[半线性波动方程的通量估计。]\label{theorem:ch12-半线性波动通量估计}假设 $u$ 是半线性波动方程\eqref{align:ch12-半线性波动} 的光滑解。
\begin{enumerate}[(i)]
\item 对每个点 $(x_0,t_0)\in\mathbb{R}^n\times(0,\infty)$，有恒等式
\begin{align}
\frac{1}{\sqrt{2}}\int_{\Gamma(x_0,t_0)}\frac{1}{2}|u_t\nu-Du|^2+F(u)\,\mathrm{d}S=e(0),
\label{align:ch12-通量恒等式}
\end{align}
其中 $\nu\coloneqq\dfrac{x-x_0}{|x-x_0|}$，
\begin{align*}
e(t)\coloneqq\int_{B(x_0,t_0-t)}\frac{1}{2}(u_t^2+|Du|^2)+F(u)\,\mathrm{d}x \quad (0\leqslant t\leqslant t_0).
\end{align*}
\item 若 $F\geqslant 0$ 且 $u(\cdot,0),u_t(\cdot,0)\equiv 0$ 在 $B(x_0,t_0)$ 内，则 $u\equiv 0$ 在锥 $K(x_0,t_0)$ 内。
\end{enumerate}
\end{theorem}

\eqref{align:ch12-通量恒等式} 左端的表达式是通过弯曲曲面 $\Gamma(x_0,t_0)$ 的能量通量。

\begin{proof}
1. 计算
\begin{align}
\dot{e}(t)&=\int_{B(x_0,t_0-t)}u_tu_{tt}+Du\cdot Du_t+f(u)u_t\,\mathrm{d}x-\int_{\partial B(x_0,t_0-t)}\frac{1}{2}(u_t^2+|Du|^2)+F(u)\,\mathrm{d}S\nonumber\\
&=\int_{\partial B(x_0,t_0-t)}\frac{\partial u}{\partial\nu}u_t-\frac{1}{2}(u_t^2+|Du|^2)-F(u)\,\mathrm{d}S\nonumber\\
&=-\int_{\partial B(x_0,t_0-t)}\frac{1}{2}|u_t\nu-Du|^2+F(u)\,\mathrm{d}S,
\label{align:ch12-e导数}
\end{align}
因为 $|u_t\nu-Du|^2=u_t^2-2u_t\dfrac{\partial u}{\partial\nu}+|Du|^2$。在 $0$ 和 $t_0$ 之间对时间积分得到\eqref{align:ch12-通量恒等式}。注意因子 $\dfrac{1}{\sqrt{2}}$ 在切换到 $\Gamma(x_0,t_0)$ 上的积分时出现，因为此曲面在 $B(x_0,t_0)\times\{t=0\}$ 上方以恒定角度 $\pi/4$ 倾斜。

2. 若 $u(\cdot,0),u_t(\cdot,0)\equiv 0$ 在 $B(x_0,t_0)$ 上，则 $e(0)=0$，因 $F(0)=0$。由于 $F\geqslant 0$，由\eqref{align:ch12-e导数} 得 $e\equiv 0$。我们推出 $u_t,Du\equiv 0$，从而 $u\equiv 0$，在锥 $K(x_0,t_0)$ 内。
\end{proof}

虽然温和拟线性波动方程\eqref{align:ch12-拟线性波动} 一般不具有守恒能量，我们仍可改编上述证明来证明有限传播速度。

\begin{theorem}[依赖域。]\label{theorem:ch12-依赖域}假设 $f(0,0,0)=0$，且 $u$ 是拟线性波动方程\eqref{align:ch12-拟线性波动} 的光滑解。若 $u(\cdot,0),u_t(\cdot,0)\equiv 0$ 在 $B(x_0,t_0)$ 内，则 $u\equiv 0$ 在锥 $K(x_0,t_0)$ 内。

因此源自 $B(x_0,t_0)$ 外部的任何扰动不影响 $K(x_0,t_0)$ 内的解。非零初值的影响以至多 $1$ 的速度传播。
\end{theorem}

\begin{proof}
定义
\begin{align*}
e(t)\coloneqq\frac{1}{2}\int_{B(x_0,t_0-t)}u_t^2+|Du|^2+u^2\,\mathrm{d}x \quad (0\leqslant t\leqslant t_0).
\end{align*}
则
\begin{align*}
\dot{e}(t)&=\int_{B(x_0,t_0-t)}u_tu_{tt}+Du\cdot Du_t+uu_t\,\mathrm{d}x-\frac{1}{2}\int_{\partial B(x_0,t_0-t)}u_t^2+|Du|^2+u^2\,\mathrm{d}S\\
&=\int_{B(x_0,t_0-t)}u_t(u_{tt}-\Delta u+u)\,\mathrm{d}x+\int_{\partial B(x_0,t_0-t)}\frac{\partial u}{\partial\nu}u_t\,\mathrm{d}S-\frac{1}{2}\int_{\partial B(x_0,t_0-t)}u_t^2+|Du|^2+u^2\,\mathrm{d}S\\
&\leqslant\int_{B(x_0,t_0-t)}u_t(-f(Du,u_t,u)+u)\,\mathrm{d}x.
\end{align*}
由于 $f(0,0,0)=0$ 且 $u$ 光滑，
\begin{align*}
|f(Du,u_t,u)|\leqslant C(|Du|+|u_t|+|u|)
\end{align*}
对某常数 $C$ 依赖于 $\|Du,u_t,u\|_{L^\infty}$。我们得出
\begin{align*}
\dot{e}(t)\leqslant C\int_{B(x_0,t_0-t)}u_t^2+|Du|^2+u^2\,\mathrm{d}x=Ce(t).
\end{align*}
由于 $e(0)=0$，Gronwall 不等式（\S B.2）蕴含 $e\equiv 0$。因此 $u\equiv 0$ 在锥 $K(x_0,t_0)$ 内。
\end{proof}

\end{document}